\documentclass[11pt,letterpaper]{article}
\usepackage[T1]{fontenc}
\usepackage{lmodern,microtype,amsmath,amssymb,amsthm,mathtools}
\usepackage[margin=.96in]{geometry}
\usepackage{xcolor}
\definecolor{inkblue}{HTML}{164567}
\usepackage[colorlinks=true,citecolor=inkblue,linkcolor=inkblue,urlcolor=inkblue]{hyperref}
\hypersetup{pdftitle={Which four-point prime variances determine the multiplicities of zeta zeros?}}
\newtheorem{theorem}{Theorem}[section]
\newtheorem{proposition}[theorem]{Proposition}
\newtheorem{corollary}[theorem]{Corollary}
\newtheorem{lemma}[theorem]{Lemma}
\newcommand{\R}{\mathbb R}
\newcommand{\Z}{\mathbb Z}
\newcommand{\dd}{\,\mathrm d}
\newcommand{\ee}{\mathrm e}
\newcommand{\ii}{\mathrm i}
\newcommand{\abs}[1]{\lvert#1\rvert}
\newcommand{\norm}[1]{\lVert#1\rVert}
\newcommand{\sym}{\mathcal D}
\newcommand{\BB}{\mathcal B}
\setlength{\emergencystretch}{1.5em}
\title{Which four-point prime variances determine\\the multiplicities of zeta zeros?}
\author{}
\date{November 27, 2026}
\begin{document}
\maketitle
\begin{abstract}
A smoothed prime-counting error has a long-time variance whose spectral weights are the squared multiplicities of zeta zeros. We classify all four-point filters on equally spaced observations for which the small-lag variance determines the limiting normalized count for every positive spectrum. Write the filter polynomial as $(1-z)^2(a+bz)$ and put $t=ab/(a^2+b^2)$. Every shape outside $(1/13,2/11)$ determines the count. Inside this interval a countable dense set of shapes misses one logarithmic frequency per shape. We characterize the set exactly. Even at an exceptional shape, every subsequential limit of the rescaled spectral measure has a density with a single sinusoidal term in the logarithm of height. This bounds the possible oscillation in the zeta multiplicity count. An integer countermodel with multiplicities one and two reproduces both Riemann--von Mangoldt counting terms. The measure theorems are unconditional. Using the repository's RH input and local multiplicity estimate, the half-width of possible limiting zeta counts is at most $0.025406$ whenever any one four-point variance has a limit.
\end{abstract}
\section{The arithmetic question}
Set $S(x)=\sum_{n\le x}\Lambda(n)(n^{-1}-x^{-1})$, where $\Lambda$ is the von Mangoldt function, and write $\gamma_{\!E}$ for Euler's constant. The ProveIt article~\cite{ProveIt} studies
\[
 E(y)=\ee^{y/2}\{S(\ee^y)-y+\gamma_{\!E}+1\}.
\]
Using its input for the Riemann zeta function, it obtains an absolutely convergent cosine series with a decaying error,
\[
 E(y)=2\sum_{\gamma>0}\frac{m_\gamma\cos(\gamma y)}{1/4+\gamma^2}+O((1+y)\ee^{-y/2}).
\]
Here $\gamma$ ranges over distinct positive ordinates and $m_\gamma$ is the order at $1/2+\ii\gamma$. For distinct real nodes $u_j$ and a nonzero real vector $c=(c_j)$ satisfying $\sum c_j=\sum c_j u_j=0$, let $p_c(x)=\sum_jc_j\ee^{\ii u_jx}$. Orthogonality in long-time mean gives
\begin{equation}\label{eq:variance}
 V_c(h)=\lim_{Y\to\infty}\frac1Y\int_0^Y\Big|\sum_jc_jE(y+hu_j)\Big|^2\dd y
   =2\sum_{\gamma>0}\frac{m_\gamma^2\abs{p_c(h\gamma)}^2}{(1/4+\gamma^2)^2}.
\end{equation}
We can translate all nodes to make them nonnegative. The $Y$ limit precedes the limit of small positive $h$. Write
\[
 \tau(h)=h^3\log(1/h),\quad M_2(T)=\sum_{0<\gamma\le T}m_\gamma^2,\quad
 \kappa(T)=\frac{2\pi M_2(T)}{T\log T}.
\]
The repository proves that one three-point variance recovers a limiting $\kappa$. Its five-point example has a convergent normalized variance for a positive measure whose normalized count oscillates. It asks which other individual stencils characterize arbitrary finite limits. We give the complete answer for four consecutive observations and a criterion for any finite set of real nodes. For the actual zeta function the results give implications and quantitative restrictions on possible limits. Whether $\kappa(T)$ itself converges remains open in this argument.

\section{The four-point classification}\label{sec:class}
On nodes $0,1,2,3$ the cancellation conditions are equivalent to
\[
 P(z)=\sum_{j=0}^3 c_jz^j=(1-z)^2(a+bz),\qquad (c_j)=(a,b-2a,a-2b,b).
\]
Reflection interchanges $a$ and $b$. Scaling and reflection preserve $t=ab/(a^2+b^2)\in[-1/2,1/2]$. Conversely $t$ determines the normalized squared response
\begin{equation}\label{eq:affine}
 \frac{\abs{P(\ee^{\ii x})}^2}{a^2+b^2}=\abs{1-\ee^{\ii x}}^4(1+2t\cos x).
\end{equation}
Call a shape \emph{determining} if its normalized variance determines the normalized counting limit for every positive measure in Theorem~\ref{thm:bank}. This implies the same recovery for~\eqref{eq:variance}.

\begin{theorem}[Four consecutive observations]\label{thm:four}
Put $I=(1/13,2/11)$ and $\alpha=\log 3/\log 2$. Every $t\notin I$ is determining. For $t\in I$ put $A=4-7t$, $B=8-32t$, $C=27t$. There are unique $u\in(0,\pi)$ and $v\in(-\pi,0)$ satisfying $B\ee^{\ii u}+C\ee^{\ii v}=A$.
The shape fails to be determining precisely when
\begin{equation}\label{eq:exception}
 (v-\alpha u)/(2\pi)\ \in\ \Z+\alpha\Z .
\end{equation}
Such shapes form a countable dense subset of $I$. Each one misses exactly one pair of frequencies $\{\pm\omega\}$. For every determining shape and finite $d\ge0$,
\begin{equation}\label{eq:arith}
 V_c(h)/\tau(h)\longrightarrow \BB_c d
 \quad\Longleftrightarrow\quad \kappa(T)\longrightarrow d,\qquad
 \BB_c=(a^2+b^2)(4+2t)/6 .
\end{equation}
On the left any decreasing mesh $h_n\to0$ with $h_{n+1}/h_n\to1$ suffices.
\end{theorem}
\begin{proof}
Theorem~\ref{thm:bank} assigns the stencil the symbol $\sum_{i<j}c_ic_j(u_j-u_i)^{3+\ii\xi}$. On four consecutive nodes, division by $a^2+b^2$ gives
\begin{equation}\label{eq:symbol}
 D_t(\xi)=7t-4+8(1-4t)\ee^{\ii\xi\log2}+27t\ee^{\ii\xi\log3}.
\end{equation}
The three lag autocorrelations, in increasing order, are $7t-4$, $1-4t$, $t$ after the same division. The normalizing constant is $(a^2+b^2)D_t(0)/6$. The zeta denominator is handled after Theorem~\ref{thm:bank}.

For $t\le0$, the magnitude of the second term exceeds the sum of the other two by $4+2t$. For $t\ge1/4$ the third exceeds the other two by $4+2t$. For $0<t<1/4$, the magnitudes $A,B,C$ are positive and satisfy
\[
 B+C-A=4+2t,\qquad B-C-A=4-52t,\qquad C-B-A=66t-12 .
\]
They form a strict triangle precisely on $I$. At its endpoints a zero requires the two phases to be $(1,-1)$ or $(-1,1)$, impossible since $\alpha$ is irrational by unique factorization.

In $I$ the phases must form one of the two conjugate configurations. For the stated orientation, lifting gives $\xi\log2=u+2\pi m$ and $\xi\log3=v+2\pi n$, proving~\eqref{eq:exception} in both directions. Irrationality makes $(m,n)$ unique. Conjugation gives the opposite frequency. The cosine rule makes $u,v$ real analytic on $I$. At the left endpoint $(u,v)\to(0,-\pi)$ and at the right $(u,v)\to(\pi,0)$. Hence $v-\alpha u$ is nonconstant on every open subinterval. Its level sets are discrete in $I$. The group $\Z+\alpha\Z$ is countable and dense, so its preimage in~\eqref{eq:exception} is countable and dense.
\end{proof}

There is also an exact test for a gap in the symbol. The irrational flow $(2^{\ii\xi},3^{\ii\xi})$ is dense on the two-torus. Outside $\overline I$, strict domination gives $\inf_\xi\abs{D_t(\xi)}>0$. Throughout $\overline I$, including its two determining endpoints, this infimum is zero because the torus contains a cancelling triangle. Almost every shape there recovers the limiting count despite arbitrarily small symbol values.

\paragraph{An exceptional shape.}
Let $F(t)=(v-\alpha u)/(2\pi)$. Choose the group element $1-\alpha$. It lies strictly between the endpoint values $-\alpha/2$ and $-1/2$, since $2^3<3^2$ and $3<2^2$. The intermediate value theorem gives an exceptional shape. For one solution $t_*$ the positive frequency is $\omega_*=(2\pi-u)/\log2$. The accompanying interval calculation sharpens the location to
\[
 \begin{split}
 0.16788340058442558&<t_*<0.16788340058442560,\\
 6.23399540333078&<\omega_*<6.23399540333080 .
 \end{split}
\]
It encloses opposite signs of $F(t)-(1-\alpha)$ at the two rational endpoints. One actual stencil is obtained by taking $a=1$ and $b=2t_* /(1+\sqrt{1-4t_*^2})$. All four entries are nonzero. Every affine-annihilating three-node stencil determines the count: for nodes $0,r,r+s$ its weights are proportional to $(s,-r-s,r)$. The longest lag term strictly dominates the other two because $(r+s)^2>r^2+s^2$. Thus four is the least number of observations at which this universal recovery can fail.

\paragraph{Two shapes suffice together.}
Let $W_t(h)=V_c(h)/(a^2+b^2)$. Equation~\eqref{eq:affine} holds pointwise under the spectral integral, so distinct shapes recover the ordinary second-difference variance at each $h$ by
\[
 W_0(h)=\frac{t_2W_{t_1}(h)-t_1W_{t_2}(h)}{t_2-t_1}.
\]
Existence of both finite normalized limits implies existence of $\lim\kappa(T)$ by~\eqref{eq:arith} at $t=0$ and fixes both original limits. The coefficients quantify how additive errors are amplified when the shapes are close.
\section{A finite test for arbitrary stencils}\label{sec:general}
Let $0\le u_1<\cdots<u_m$ and let $c$ satisfy the two cancellations of Section~1. Define, at $x=0$ by continuous extension,
\[
 K_c(x)=\frac{\abs{p_c(x)}^2}{x^4},\quad
 \BB_c=\frac1\pi\int_0^\infty K_c(x)\dd x,\quad
 \sym_c(\xi)=\sum_{i<j}c_ic_j(u_j-u_i)^{3+\ii\xi}.
\]
For a positive locally finite Borel measure $A$ on $[0,\infty)$ with $\int_1^\infty x^{-4}\dd A(x)<\infty$, write $A(T)=A([0,T])$ and $U_c(h)=h^4\int K_c(hx)\dd A(x)$. An atom at zero contributes $O(h^4)$.

\begin{theorem}[Finite families of filters]\label{thm:bank}
For a nonempty finite family of nonzero stencils the following property holds exactly when the functions $\sym_c$ have no common real zero. For all such $A$ and finite $d\ge0$,
\[
 \frac{\pi A(T)}{T\log T}\longrightarrow d
 \quad\Longleftrightarrow\quad
 \frac{U_c(h)}{\tau(h)}\longrightarrow d\BB_c\quad\text{for every stencil in the family.}
\]
For the converse the variances need only be sampled on a mesh with $h_{n+1}/h_n\to1$. Each $\BB_c=\sym_c(0)/6>0$.
\end{theorem}
\begin{proof}
The cancellations imply $p_c(x)=O(\min(x^2,1))$. For $F(y)=\cos y-1+y^2/2$, expanding the squared modulus gives
\begin{equation}\label{eq:comp}
 \abs{p_c(y)}^2=2\sum_{i<j}c_ic_jF((u_j-u_i)y).
\end{equation}
The constant and quadratic corrections have the required sums since $\sum c_j=\sum u_jc_j=0$. Twice integrating by parts gives
$\int_0^\infty F(x)x^{-4}\dd x=\frac16\int_0^\infty(1-\cos x)x^{-2}\dd x=\pi/12$ by the Dirichlet integral. Equation~\eqref{eq:comp} proves the formula for $\BB_c$. Positivity follows from $K_c\ge0$ and linear independence of exponentials at distinct nodes.

Consider $\mu_T(B)=\pi A(TB)/(T\log T)$ for $T>2$. A counting limit $d$ implies vague convergence to $d\dd x$ by evaluation on intervals. The bound $A(T)\ll(1+T)\log(2+T)$ gives uniformly small tails for the $K_c$ integrals. For $R\ge1$ the mass of $[R,2R]$ is $O(R(1+\log(2R)))$, uniformly for $T>2$. Dyadic summation gives $\int_R^\infty x^{-4}\dd\mu_T(x)=O(R^{-3}(1+\log(2R)))$. With $h=1/T$ we get the desired limit since $K_c(x)=O(x^{-4})$ at infinity.

Conversely choose one stencil. For a small $\delta>0$, $\abs{p_c}$ is bounded below on $[\delta/3,2\delta]$ by analyticity and nontriviality. Take mesh indices with $h_n\sim\delta/T$. Then $h_nx$ stays in that interval for $T/2\le x\le T$. As $U_c(h_n)=O(\tau(h_n))$, this gives $A([T/2,T])\ll T\log T$. Dyadic summation gives the cumulative growth bound above, including when the stencil has higher-order cancellation at the origin.

Take a vaguely convergent subsequence of $\mu_T$, with limit $\nu$. Such subsequences exist by local uniform boundedness, and $\nu([0,R])\ll R$ for every $R>0$ by the growth bound and passage to the limit. For each $s>0$, choose mesh indices with $Th_n\to s$. The continuous kernels converge on compacts, the tail bounds hold uniformly for $Th_n$ near $s$, and $\log(1/h_n)/\log T\to1$. Since
\[
 \frac{U_c(s/T)}{T^{-3}\log T}
     =\frac1\pi\int_0^\infty \abs{p_c(sx)}^2 x^{-4}\dd\mu_T(x),
\]
the signed measure $\sigma=\nu-d\dd x$ satisfies $\int\abs{p_c(sx)}^2x^{-4}\dd\sigma(x)=0$ for every $s>0$ and each stencil.

Put $R(s)=\int F(sx)x^{-4}\dd\sigma(x)$. Since $\abs{\sigma}([0,X])\ll X$ and $\abs{F(y)}\ll\min(y^4,y^2)$, splitting at $1/\abs s$ and summing the tail gives $R(s)=O(\abs{s}^3)$. This defines a continuous even function. Equation~\eqref{eq:comp} gives $\sum_{i<j}c_ic_jR((u_j-u_i)s)=0$. In the variable $w=\log s$ this is a finite translation equation for the bounded function $H(w)=\ee^{-3w}R(\ee^w)$. With $\widehat f(\xi)=\int\ee^{-\ii w\xi}f(w)\dd w$, its distributional transform satisfies $\sym_c(\xi)\widehat H(\xi)=0$. If the symbols have no common real zero, test against compactly supported smooth functions divided by $\sum_c\abs{\sym_c}^2$, and sum the equations multiplied by $\overline{\sym_c}$. Thus $\widehat H=0$ and $R=0$.

Finally, $R^{(4)}$ is the Fourier transform of the symmetrization of $\sigma$ (give each of the two reflections weight $1/2$). Indeed $\partial_s^4 F(sx)=x^4\cos(sx)$, and truncation followed by testing against Schwartz functions justifies the identity. The growth bounds make the measures tempered and justify passage to the limit. Fourier injectivity gives $\sigma=0$. Every subsequential limit is therefore $d\dd x$, and evaluation on $[0,1]$ proves the counting limit.

For necessity let $\omega$ be a common zero. It is nonzero since $\BB_c>0$. For real $\xi$ substitution in~\eqref{eq:comp} gives
\begin{equation}\label{eq:mellin}
 \int_0^\infty K_c(y)y^{\ii\xi}\dd y
 =2\sym_c(-\xi)\int_0^\infty F(y)y^{-4+\ii\xi}\dd y .
\end{equation}
The integrals converge absolutely. On $x\ge1$ take $\dd A(x)=\pi^{-1}\log x\,[1+\epsilon\cos(\omega\log x)]\dd x$, $0<\epsilon<1$. Its normalized count is $1+\epsilon\Re(T^{\ii\omega}/(1+\ii\omega))+o(1)$. Yet substitution $y=hx$ gives $U_c(h)/\tau(h)\to\BB_c$. The oscillatory coefficient in~\eqref{eq:mellin} vanishes, $K_c(y)\abs{\log y}$ is integrable, and the integral on $0<y<h$ lost by the cutoff tends to zero.
\end{proof}

In fact, for a family with no common zero, mere existence of all finite normalized variance limits implies their proportionality to $\BB_c$. To see this, repeat the compactness argument with the individual limits. Each tangent $\nu$ gives the same value as its dilation $\nu(s\,\cdot)/s$ in every scaled kernel integral: both are tangents of the same family. Applying the injectivity argument to their difference gives $\nu(s\,\cdot)=s\nu$ for every $s>0$. Thus $\nu$ has constant density on the half-line, and any one measured limit fixes that constant for all tangents.

For zeta use $\dd A=2\sum_{\gamma>0}m_\gamma^2\delta_\gamma$. The repository's local multiplicity bound gives $M_2(T)\ll T\log(2+T)$, hence $\int(1+x^2)^{-1}\dd A(x)<\infty$. For $x>0$ one has $0\le1-x^4/(1/4+x^2)^2\ll(1+x^2)^{-1}$. The bound $\abs{p_c(hx)}^2\ll h^4x^4$ shows that replacing $x^{-4}$ by $(1/4+x^2)^{-2}$ costs $O(h^4)=o(\tau(h))$. This proves the arithmetic implication in Theorem~\ref{thm:four}.

\section{What an exceptional variance still determines}\label{sec:tangent}
\begin{theorem}[One missing frequency]\label{thm:tangent}
Suppose $t$ is exceptional, with positive frequency $\omega$, and $U_c(h)/\tau(h)\to d\BB_c$ on a mesh as above. Every subsequential vague limit of $\mu_T$ has the form
\begin{equation}\label{eq:tan}
 [d+\Re(Cx^{\ii\omega})]\dd x,\qquad C\in\mathbb C,\qquad \abs C\le d.
\end{equation}
There is some $r\in[0,d]$ such that the full set of subsequential limits of $\pi A(T)/(T\log T)$ is
\[
 \left[d-\frac r{\sqrt{1+\omega^2}},\,d+\frac r{\sqrt{1+\omega^2}}\right].
\]
Put $L=\exp(\pi/\omega)$ and $a(T)=\pi A(T)/(T\log T)$. There is also a pointwise-in-height cancellation,
\[
 a(T)+a(LT)\ \longrightarrow\ 2d .
\]
For $A=2\sum m_\gamma^2\delta_\gamma$ the repository's local bounds sharpen this to
\begin{equation}\label{eq:cap}
 r\le\min(d-1,\,1.3208-d).
\end{equation}
In particular a limiting variance at any four-point shape fixes the midpoint of the lower and upper limiting normalized zeta counts. If $d=1$ or $d=1.3208$, the whole count converges. At the lower endpoint, the repository identifies $\kappa(T)\to1$ with density-one simplicity.
\end{theorem}
\begin{proof}
The compactness and translation argument of Theorem~\ref{thm:bank} applies to every limiting $\nu$. For $\sigma=\nu-d\dd x$, the bounded $H$ has Fourier transform supported on $\{\pm\omega\}$. A distribution supported on finitely many points is a sum of derivatives of delta masses. Thus $H$ is a sum of exponentials at these frequencies multiplied by polynomials. Boundedness forces the polynomials to be constant, since the top coefficients form a nonzero periodic trigonometric polynomial bounded away from zero along some sequence of period translates. Reality gives conjugate coefficients.

For the measures $\Re(Cx^{\ii\omega})\dd x$ the transform $\sigma\mapsto H$ is $\Re(CJ\ee^{-\ii\omega w})$, where $J=\int_0^\infty F(x)x^{-4+\ii\omega}\dd x$. It is nonzero. If $J=0$, then $R$ for $\cos(\omega\log x)\dd x$ would vanish identically, contradicting the injectivity of $\sigma\mapsto R$ established using its fourth derivative. Hence this two-dimensional real family realizes every possible $H$. Injectivity gives~\eqref{eq:tan}, and positivity gives $\abs C\le d$.

The family of subsequential limits is compact for vague convergence. Its parameter $C$ is continuous because two compactly supported test integrals distinguish the independent functions $\cos(\omega\log x)$ and $\sin(\omega\log x)$. Let $r$ be the largest modulus attained. Replacing $T$ by $sT$ replaces $\nu(B)$ by $\nu(sB)/s$, rotating $C$ by $s^{\ii\omega}$. The full circle of radius $r$ occurs. Evaluation on $[0,1]$ gives $d+\Re(C/(1+\ii\omega))$ and proves the interval statement. For $s=L$ the rotation negates $C$, so the sum of the two counts tends to $2d$ along every tangent subsequence. Compactness gives convergence along all heights.

For zeta, the repository's local squared-multiplicity lemma~\cite{ProveIt} bounds every tangent between $\dd x$ and $1.3208\dd x$ on each fixed positive interval. These interval inequalities bound the density almost everywhere and hence everywhere on $(0,\infty)$ by continuity. Its extreme values are $d\pm\abs C$, giving~\eqref{eq:cap}.
\end{proof}
\begin{proposition}[The first missing frequency]\label{prop:firstres}
The least positive frequency missed by any four-point shape is the number
$\omega_*$ enclosed in Section~\ref{sec:class}. Thus a convergent four-point
variance restricts the half-width of the limiting normalized zeta counts to
\[
 \frac{0.1604}{\sqrt{1+\omega_*^2}}=0.025405106245\ldots .
\]
This constant is sharp for spectra satisfying the two local bounds used in~\eqref{eq:cap}.
\end{proposition}
\begin{proof}
In the triangle interval, $A=B/2+C/3$, with $B,C>0$. A zero therefore requires
$2^{\ii\omega}-1/2$ and $3^{\ii\omega}-1/3$ to point in opposite directions.
Write $L=2\pi/\log3$ and $U=3\pi/(2\log2)$. For $0<\omega\le L$ the arguments
$\omega\log2$ and $\omega\log3$ bound an arc of length at most
$2\pi\log(3/2)/\log3<\pi$. The positive combination
$B2^{\ii\omega}+C3^{\ii\omega}$ points strictly inside that arc, whose interior excludes
the positive real axis. Hence this range contains no missing frequency.

Let $\theta_p(\omega)$ be the continuous argument of
$p^{\ii\omega}-1/p$ starting at zero. With $\phi=\omega\log p$,
\[
 \theta'_p=(\log p)\frac{1-p^{-1}\cos\phi}{1-2p^{-1}\cos\phi+p^{-2}}.
\]
On $[L,U]$ the angle $\omega\log2$ lies between $\pi$ and $3\pi/2$,
so $\theta'_2\le\log2$. The same formula gives
$\theta'_3\ge(3/4)\log3>\log2$. Thus $\theta_3-\theta_2$ strictly increases.
At $L$ its value lies in $(0,\pi)$ since $\theta_3(L)=2\pi$ and
$\pi<\theta_2(L)<2\pi$. At $U$ we have $\theta_2(U)=\pi+\arctan 2$,
whereas $\theta_3(U)=2\pi+\arg(\ee^{\ii\psi}-1/3)$ with
$\psi=(3\alpha/2-2)\pi\in(\pi/3,\pi/2)$, using $3^9>2^{14}$
and $3^3<2^5$. Here $\sin\psi>2(\cos\psi-1/3)$, so
$\pi<\theta_3(U)-\theta_2(U)<2\pi$. The unique crossing of $\pi$
is the first possible opposition. The enclosed pair $(t_*,\omega_*)$
already realizes a zero in $(L,U)$, hence gives this crossing.
Equation~\eqref{eq:cap} gives $r\le(1.3208-1)/2$;
the discrete construction in Section~\ref{sec:discrete} attains equality at $\omega_*$.
\end{proof}


\section{Integer multiplicities with the counting main terms}\label{sec:discrete}
\begin{proposition}[Discrete realization]\label{prop:discrete}
For every exceptional shape, every $0<q_0<1$, and $0<\epsilon\le\min(q_0,1-q_0)$, there is a locally finite set of positive ordinates with multiplicities $m_x\in\{1,2\}$ such that
\begin{align*}
 \sum_{x\le T}m_x&=\frac{T}{2\pi}\log\frac{T}{2\pi}-\frac{T}{2\pi}+O(1),\\
 \frac{2\pi}{T\log T}\sum_{x\le T}m_x^2
 &=1+q_0+\epsilon\Re\frac{T^{\ii\omega}}{1+\ii\omega}+o(1).
\end{align*}
The spectral variance in~\eqref{eq:variance} obeys $V_c(h)/\tau(h)\to(1+q_0)\BB_c$. Thus the factor $1/\sqrt{1+\omega^2}$ is attained with an integer spectrum whose ordinary count agrees with both classical main terms.
\end{proposition}
\begin{proof}
For $x\ge10$ put $\lambda(x)=\log(x/(2\pi))/(2\pi)$ and $q(x)=q_0+\epsilon\cos(\omega\log x)$. Define the strictly increasing unbounded functions
\[
 F_1(T)=\int_{10}^T(1-q(x))\lambda(x)\dd x,\qquad F_2(T)=\frac12\int_{10}^T q(x)\lambda(x)\dd x .
\]
Place simple points at the positive integer crossings of $F_1$ and double points at the positive crossings of $\Z+\theta$ by $F_2$. Choose $0<\theta<1$ outside the countable set of fractional parts of $F_2(F_1^{-1}(n))$ for positive integers $n$. The point sets are disjoint. The two counts differ from $F_1,F_2$ by $O(1)$ each. Thus the ordinary count is $\int_{10}^T\lambda+O(1)$ and the squared count is $\int_{10}^T(1+q)\lambda+O(1)$. Direct integration gives the asserted formulas.

The continuous squared-counting density $2(1+q)\lambda$ has variance limit $(1+q_0)\BB_c$ by~\eqref{eq:mellin}, since the constant in $\log(x/(2\pi))$ contributes $O(h^3)$. Replacing the measure by $2\sum m_x^2\delta_x$ changes its cumulative function by $O(1)$. Integration by parts against $h^4K_c(hx)$ costs $O(h^4\int_0^\infty\abs{K'_c(y)}\dd y)=O(h^4)$, as $K_c$ is smooth at zero and $K'_c=O(y^{-4})$ at infinity. The zeta denominator change also costs $O(h^4)$.
\end{proof}
Taking $q_0=\epsilon=0.1604$ retains the local limiting bounds $1$ and $1.3208$ on every fixed interval while attaining the full half-width allowed by those bounds. The constructed ordinates are a model spectrum. An Euler product realizing them would be an additional arithmetic constraint.

\section*{Relation to prior work and further questions}
The ProveIt source~\cite{ProveIt} supplies the prime explicit formula, the three-point inverse theorem, a five-point model and the local number-theoretic upper bound. Its converse for three points uses a strict contraction. Multiplicative cancellation, Mellin nonvanishing and log-periodic counterexamples occur in inverse regular variation~\cite{JMRS,DMRS}. Greenhall classified ambiguity for the Allan variance, another scale-dependent filtered variance~\cite{Greenhall}. Our contribution is the distance symbol for this prime-variance problem, its four-node classification and single missing frequency, and the resulting limits on possible zeta multiplicity counts. A comparison with these primary papers and the pinned repository found no earlier statement of this classification. Wider priority remains open.

For five equally spaced observations the symbol uses the dependent phases $2^{\ii\xi}$ and $4^{\ii\xi}$ together with $3^{\ii\xi}$. Which shapes miss multiple positive frequencies? For zeta itself, can arithmetic constraints beyond the local bounds exclude a nonconstant density in~\eqref{eq:tan}? A positive answer would remove the exceptional set from the arithmetic equivalence.

\begin{thebibliography}{9}
\bibitem{ProveIt} ProveIt Contributors, \emph{Beyond the Riemann hypothesis for Dirichlet $L$-functions: consequences, transfers and limits}, snapshot \texttt{35e2a5a63492} (2026). See the subsection on merging scales, especially \texttt{zer:lem:localM2}, \texttt{zer:thm:inversemeasure}, and \texttt{zer:thm:coalesceinverse}. \href{https://github.com/VladimirReshetnikov/ProveIt/blob/35e2a5a634924e57ac6b77bd73a28570dd1a362b/Analysis/GeneralizedRiemannHypothesis/Article/sections/zer.tex}{Pinned variance and inverse problem source}. \href{https://github.com/VladimirReshetnikov/ProveIt/blob/35e2a5a634924e57ac6b77bd73a28570dd1a362b/Analysis/GeneralizedRiemannHypothesis/Article/sections/src.tex}{Repository input and scope}.
\bibitem{JMRS} M. Jacobsen, T. Mikosch, J. Rosi\'nski and G. Samorodnitsky, \emph{Inverse problems for regular variation of linear filters, a cancellation property for $\sigma$-finite measures and identification of stable laws}, Ann. Appl. Probab. \textbf{19} (2009), 210--242. \href{https://doi.org/10.1214/08-AAP540}{doi:10.1214/08-AAP540}.
\bibitem{DMRS} E. Damek, T. Mikosch, J. Rosi\'nski and G. Samorodnitsky, \emph{General inverse problems for regular variation}, J. Appl. Probab. \textbf{51A} (2014), 229--248. \href{https://doi.org/10.1239/jap/1417528478}{doi:10.1239/jap/1417528478}.
\bibitem{Greenhall} C. A. Greenhall, \emph{Spectral ambiguity of Allan variance}, IEEE Trans. Instrum. Meas. \textbf{47} (1998), 623--627. \href{https://doi.org/10.1109/19.744312}{doi:10.1109/19.744312}.
\end{thebibliography}
\end{document}
